\documentclass[12pt]{article}
%^ Start of document
\usepackage{times}  %Makes text TimesNewRoman
\usepackage{setspace} %Allows you to set single or double space 
\usepackage{biblatex} %Imports biblatex package
\usepackage{graphicx} % Required for inserting images
\usepackage{authblk}
\usepackage[colorlinks=true, urlcolor=blue, linkcolor=blue]{hyperref}
\usepackage{subfig}
\usepackage{float}
\usepackage{tabularx}
\usepackage{multirow}
\usepackage[paperheight=11in,
   paperwidth=8.5in,
   top=20mm,
   bottom=20mm,
   left=40mm,
   right=40mm]{geometry}
\usepackage{moresize}
\usepackage{tikz}
\usepackage{xcolor}
\usepackage{physics}
\usepackage{amssymb}
\usepackage{mathrsfs}
\usepackage{amsmath}
\usepackage{ragged2e}
\usepackage{comment}
\usepackage{xcolor}
\addbibresource{seminar.bib}
\begin{document}
\singlespacing  %Sets single spacing for whole document
\title{Astronomy}

\author[1]{Zachary Tullier}
\affil[1]{Student\\Physics Department\\ University of Houston - Clear Lake \\Houston, Texas\\U.S}
\date{}
\maketitle




\section{Inspirational Philosophy for Science and Life}
    Pursue Your Passion. Follow what you genuinely love as the foundation for a fulfilling career. True success comes from engaging in work that resonates with your interests. Dream Big and Aim High. Set ambitious goals and push the boundaries of what you believe is achievable. Share Knowledge Freely. Collaboration and openness are essential for both personal and scientific growth. Knowledge should be shared, not hoarded. Break Limitations. Challenge imposed limitations, whether self-inflicted or external, and push yourself beyond perceived boundaries. Physics and astronomy increasingly benefit from insights and tools from biology, chemistry, and artificial intelligence (AI). Interdisciplinary approaches drive innovation, such as using AI to enhance spectroscopy or applying biological models to understand cosmic phenomena.


\section{The Role of Artificial Intelligence (AI) in Astronomy}
    AI enables astronomers to handle complex problems by analyzing vast datasets with speed and precision. AI facilitates discoveries, such as identifying exoplanets and detecting gravitational waves, that would be impossible through manual analysis. Autonomous robotics, including planetary rovers, now rely heavily on AI for real-time decision-making. Astronomy chatbots provide real-time answers to astronomy-related questions, democratizing knowledge for diverse audiences. Citizen science engages the public in scientific work, such as galaxy classification, while improving machine learning models. AI analyzes light curves from stars to identify planetary transits and atmospheric properties. AI enables autonomous exploration of planetary surfaces, prioritizing tasks and navigating terrain based on real-time data. AI systems analyze spectra to identify elements and molecules in celestial bodies, providing insights into their composition. It improves detection of exoplanet atmospheres by identifying subtle chemical signatures during planetary transits. AI isolates signals from noise in massive datasets, aiding in the identification of gravitational waves. Multi-step noise reduction techniques and advanced pattern recognition enhance our ability to study cosmic events like black hole mergers.


\section{Overcoming Challenges in Space Exploration}
    AI organizes, processes, and interprets astronomical data that exceeds human capacity for analysis. AI systems optimize spacecraft paths using gravity assists, reducing time and energy requirements. AI supports autonomous spacecraft operation, onboard diagnostics, and real-time problem-solving for Mars and beyond.


\section{Ethical and Practical Considerations in AI}
    Ensuring that AI-driven results can be replicated and validated is essential. Over-reliance on AI may lead to errors or nonsensical outputs, highlighting the need for human oversight. The speaker discussed potential biases in algorithms and the broader implications of AI in research and decision-making.


\section{A Path Forward}
    Solving complex scientific problems requires collaboration between physics, AI, biology, and environmental sciences. Open science and data-sharing foster innovation and collective progress. The speaker suggests:
    \begin{itemize}
        \item Explore emerging fields like gravitational wave astronomy, AI in space exploration, and exoplanet research.
        \item Cultivate skills across multiple domains, such as programming, data science, and theoretical physics.
        \item Embrace curiosity-driven research and think beyond conventional boundaries.
    \end{itemize}

\section{Conclusion}
    Embrace AI and interdisciplinary approaches to push the limits of knowledge in physics and astronomy. Solve pressing challenges in space exploration, such as data management, autonomous robotics, and sustainability. Collaborate across borders and disciplines to make groundbreaking discoveries that benefit humanity.

\end{document}
